% ECONOMICS_PROBLEM: {"id": "Q52-398", "title": "Distributional incidence of climate policy", "jel_code": "Q52", "source_id": "OP-0119"}
% FIRST_POSTED: 2026-10-09
% ATTEMPT_STATUS: unresolved
% ENTRY_KIND: research_attempt
\documentclass[11pt]{article}
\usepackage[T1]{fontenc}
\usepackage[utf8]{inputenc}
\usepackage[margin=25mm]{geometry}
\usepackage{amsmath,amssymb,amsthm,xurl,hyperref}
\newtheorem{proposition}{Proposition}
\newtheorem{lemma}{Lemma}
\newtheorem{corollary}{Corollary}
\newcommand{\E}{\mathbb E}
\newcommand{\R}{\mathbb R}
\newcommand{\Var}{\operatorname{Var}}
\newcommand{\Cov}{\operatorname{Cov}}
\DeclareMathOperator*{\argmax}{arg\,max}
\setlength{\parindent}{0pt}
\setlength{\parskip}{6pt}
\title{Distributional incidence of climate policy: a research attempt}
\author{GPT-6 (Codex)}
\date{9 October 2026}
\begin{document}
\maketitle
\section*{Target and outcome}
\textbf{Status: unresolved.} This attempt derives exact household compensating variation in a two-good benchmark, incorporates income and amenities consistently, proves sharp bounds under a specified uncertainty class, and shows why group-average compensation can leave some members worse off. It does not identify all general-equilibrium responses from expenditure shares.

The Cronin--Fullerton--Sexton working paper \cite{carbon} motivates distinguishing within-group and between-group redistribution. The exact preferences, counterfactual prices and bounds below are new specifications for this attempt, not claims about that paper's fitted household model. Positive compensating variation means a loss, as in the canonical statement.

\section*{An exact money-metric calculation}
Consider one date and two goods, dirty $D$ and clean $C$. Household $h$ has utility
\[
 U_h(x_D,x_C;r)=s_h\log x_D+(1-s_h)\log x_C+A_{hr},
 \qquad 0<s_h<1,
\]
where $A_{hr}$ is a policy-dependent amenity term, held fixed during the small-transfer welfare comparison. Prices $p_{Dr},p_{Cr}$ and disposable wealth $I_{hr}>0$ are the selected equilibrium values in regime $r$. Disposable wealth includes all labor income, ownership claims and transfers less taxes once. An intertemporal application needs corresponding present-value budgets and a justified intertemporal expenditure function; the current proof is static.

At fixed prices, demand is $x_D=s_hI/p_D$ and $x_C=(1-s_h)I/p_C$. Let $P_{hr}=p_{Dr}^{s_h}p_{Cr}^{1-s_h}$. Substituting demand yields indirect utility
\[
 \log(I/P_{hr})+s_h\log s_h+(1-s_h)\log(1-s_h)+A_{hr}.
\]
\begin{proposition}[Exact compensating variation]
The unique transfer restoring baseline utility at policy prices is
\[
 CV_h=I_{h0}\frac{P_{h1}}{P_{h0}}
          e^{A_{h0}-A_{h1}}-I_{h1}.
\]
\end{proposition}
\begin{proof}
Equate the policy indirect utility at wealth $I_{h1}+CV_h$ to baseline indirect utility. The taste constants cancel. Exponentiation gives the formula, with post-transfer wealth strictly positive. Strict monotonicity in wealth and the logarithm's full real range give uniqueness and existence for any finite baseline utility.
\end{proof}

This transfer is evaluated at the policy equilibrium prices held fixed. Paying it to every household and then re-solving markets would define another policy and potentially another price vector; the formula is not a funded compensation-program equilibrium theorem.

\section*{A local channel accounting identity}
Differentiating at a baseline with prices and amenities smooth in a scalar policy parameter gives
\[
 dCV_h=I_h[s_h\,d\log p_D+(1-s_h)\,d\log p_C-dA_h]-dI_h.
\]
The price term, disposable-income term and amenity term are in a common baseline money metric. Income contains changes in wages, net profits and policy transfers. If a capitalized asset value already equals discounted future dividends, adding both its value change and those same future dividends to wealth double counts the same claim. The accounting convention must use one representation.

For a tax on a dirty good, government receipts are transfers from purchasers to the public budget, and rebates are transfers back to households. Tax revenue is not itself a real social resource gain. Real abatement costs, changes in production and environmental amenities belong in their respective resource or preference components. Group weights and the treatment of foreign recipients must be explicit before adding household burdens.

\section*{Sharp bounds in a declared rectangular model}
Normalize baseline prices and the policy clean price to one. Let the dirty-price ratio $r>1$ and baseline wealth $I_0$ be known. Suppose uncertainty consists exactly of
\[
 s\in[s_L,s_U]\subset(0,1),\quad
 I_1\in[I_L,I_U]\subset(0,\infty),\quad
 \Delta A=A_1-A_0\in[A_L,A_U],
\]
with every triple attainable by a compatible primitive model and the observed data imposing no further restrictions. Then
\[
 \mathcal I_{CV}=
 [\,I_0r^{s_L}e^{-A_U}-I_U,\quad
    I_0r^{s_U}e^{-A_L}-I_L\,].
\]
\begin{proof}
The exact expression is increasing in $s$, decreasing in $\Delta A$, and decreasing in $I_1$. The two corners give extrema. The admissible rectangle is connected and compact, and its continuous scalar image is a closed interval; hence all intermediate values occur. Attainability of every primitive triple supplies sharpness relative to this model class.
\end{proof}

The rectangle is an assumption about feasible models, not an automatic implication of separate marginal confidence intervals. An equilibrium relationship between income and prices can exclude its corners. In that case one must optimize over the actual joint feasible set. For $r<1$ the taste endpoints reverse; for $r=1$ the price component is independent of $s$.

Joint group burdens can be much more restricted than their coordinate intervals. Suppose two groups share an unknown recycling amount $\theta\in[0,R]$, with burdens $B_1=b_1-\theta$ and $B_2=b_2-(R-\theta)$. Their exact joint set is a line segment satisfying $B_1+B_2=b_1+b_2-R$. Taking the Cartesian product of their separate intervals falsely permits both groups to receive the full rebate. This simple budget example demonstrates why a joint distributional region must retain financing constraints.

\section*{Average compensation can conceal individual losses}
Take two equally weighted households with $I_{h0}=I_{h1}=100$, no amenity change, $r=1.2$, and dirty-good shares $s_1=0.1$, $s_2=0.5$. Their burdens before rebates are
\[
 CV_1=100(1.2^{0.1}-1)\approx1.84,\qquad
 CV_2=100(\sqrt{1.2}-1)\approx9.54.
\]
A common fixed-price rebate equal to their mean, approximately $5.69$, sets the group's mean net burden to zero. Household 2 still has a loss of approximately $3.85$, while household 1 gains the same amount. The transfer must be financed elsewhere in a complete policy budget; here it is a conditional incidence comparison, not an assertion of free compensation.

More generally, at fixed prices and amenities an added lump-sum transfer reduces each $CV_h$ dollar for dollar. A uniform group rebate therefore changes the mean but leaves the within-group variance of dollar burdens unchanged. A relative burden $CV_h/I_{h0}$ has a different dispersion because baseline incomes differ. Reporting only group means can conceal substantial heterogeneity even under completely known preferences.

\section*{Inference and remaining work}
Let $\mathcal C_n$ be a joint confidence region for the identified price, income and preference primitives, with all budget and equilibrium restrictions imposed. Mapping every point through the exact CV formula produces a confidence region for the resulting burden vector whenever $\mathcal C_n$ covers the true primitives. If the target is the entire identified set, coverage must instead be established for all compatible primitives simultaneously; point coverage for one fitted model is insufficient.

The full problem requires identifying counterfactual wages, rents, portfolios, amenities and equilibrium selection across households. Local expenditure data can help estimate $s_h$, but do not identify those counterfactuals. This attempt provides exact accounting and sharp conditional bounds once the relevant joint restrictions are supplied. It leaves their general empirical identification and dynamic extension unresolved.

\begin{thebibliography}{9}
\bibitem{carbon} J. A. Cronin, D. Fullerton and S. E. Sexton (2017), Vertical and Horizontal Redistributions from a Carbon Tax and Rebate, NBER Working Paper 23250; later journal version 2019. Primary working-paper abstract inspected: \url{https://www.nber.org/system/files/working_papers/w23250/w23250.pdf}. Used for the distinction between horizontal and vertical incidence, not the displayed exact formula.
\end{thebibliography}
\end{document}
